\documentclass[12pt]{article}
\usepackage[T1]{fontenc}
\usepackage[utf8]{inputenc}
\usepackage{lmodern}
\usepackage{textcomp}
\usepackage{enumitem}
\usepackage{geometry}
\geometry{margin=1in}
\begin{document}
\begin{center}
Answer Key - Sociology - Undergraduate Level Assessment 17 \
\small (Answers are ordered exactly as in the question paper.)
\end{center}

\section*{Multiple Choice}
\begin{enumerate}[label=\arabic*.]
\item \textbf{Answer:} D
\\ \textit{Options:} A) this pattern of speech made them the target of bullying; B) they referred to explicit, context independent meanings; C) they prevented children from communicating outside of their peer groups; D) they involved short, simple sentences with a small vocabulary
\\ \textit{Note:} Bernstein's theory suggests that 'restricted codes' limit the complexity and richness of language used by pupils, thereby hindering their ability to express ideas and engage in academic discourse effectively, which can disadvantage them in an educational setting.
\item \textbf{Answer:} B
\\ \textit{Options:} A) repress it, by promoting only the interests of elite groups; B) revive it, by reaffirming a commitment to freedom of speech; C) reproduce it, by emphasizing face-to-face contact with peer groups; D) replace it with a superior form of communication
\\ \textit{Note:} The Internet has revived the public sphere by providing a platform for diverse voices and facilitating open discourse, thereby reaffirming the fundamental democratic principle of freedom of speech.
\item \textbf{Answer:} D
\\ \textit{Options:} A) the product they were making; B) the production process; C) each other and humanity in general; D) all of the above
\\ \textit{Note:} Marx argued that capitalism alienates workers from the products of their labor, the labor process itself, their fellow workers, and their own human potential, leading to a comprehensive sense of disconnection in all aspects of their work and existence.
\item \textbf{Answer:} A
\\ \textit{Options:} A) Functionalism; B) Conflict theory; C) Symbolic interactionism; D) Role theory
\\ \textit{Note:} Functionalism is the correct answer because it emphasizes that society is made up of various interdependent and coordinated parts, each serving a specific function to maintain the stability and order of the whole system.
\item \textbf{Answer:} A
\\ \textit{Options:} A) decreased rates of prostitution and sex tourism; B) developing countries can depend on it as a crucial source of income; C) the exploitation of cheap, unregulated labour in poor countries; D) we have become more aware of 'other' societies and ways of living
\\ \textit{Note:} The correct answer is "decreased rates of prostitution and sex tourism" because global tourism often exacerbates these issues by increasing demand for sexual services in destination areas, rather than reducing them.
\end{enumerate}

\section*{True / False}
\begin{enumerate}[label=\arabic*.]
\setcounter{enumi}{5}
\item \textbf{Answer:} True
\\ \textit{Options:} A) True; B) False
\\ \textit{Note:} The statement is true because sociologists recognize that their interpretations of social phenomena are influenced by their own beliefs and values, which inevitably shapes the kind of knowledge they produce and consider relevant.
\item \textbf{Answer:} False
\\ \textit{Options:} A) True; B) False
\\ \textit{Note:} Secularization refers to the process by which religious institutions, practices, and beliefs lose their social significance, leading to a decline in religious participation and belief in supernatural forces, which makes the statement false.
\item \textbf{Answer:} True
\\ \textit{Options:} A) True; B) False
\\ \textit{Note:} An open society is characterized by social mobility, which enables individuals to improve their social status through education, employment opportunities, and personal achievements, thus allowing movement between different levels of the social hierarchy.
\item \textbf{Answer:} True
\\ \textit{Options:} A) True; B) False
\\ \textit{Note:} The statement is true because industrial capitalism transformed labor patterns, leading to designated leisure time that allowed for organized recreational activities and the development of civic culture, reflecting the social changes of the era.
\item \textbf{Answer:} True
\\ \textit{Options:} A) True; B) False
\\ \textit{Note:} Marxist feminists assert that women's unpaid domestic labor supports and sustains the capitalist economy by allowing men to participate in wage labor, thereby perpetuating both gender inequality and the capitalist system.
\end{enumerate}

\section*{Long Form Response}
\begin{enumerate}[label=\arabic*.]
\setcounter{enumi}{10}
\item \textbf{Answer:} Current research findings indicate that the unequal division of housework and childcare remains a significant issue among married couples, particularly when both partners are employed. Despite the increasing participation of women in the workforce, many men do not share the responsibilities of domestic labor equitably with their working wives. This disparity can perpetuate traditional gender roles and reinforce the notion that women are primarily responsible for home management, even when they are also contributing financially. The implications of this imbalance extend beyond the household, impacting women's career advancement, mental health, and overall societal perceptions of gender equality in both domestic and professional settings. Addressing these inequalities is essential for fostering a more equitable workplace and home environment.
\\ \textit{Note:} correct because it highlights the persistent gender inequality in domestic labor despite women's increased workforce participation, thereby illustrating how traditional gender roles are reinforced within the context of married couples.
\item \textbf{Answer:} Functionalist theory in sociology posits that society functions as a complex system whose components work together to promote stability and order. Key components of this theory include social institutions such as family, education, religion, and government, each contributing to the overall equilibrium by fulfilling specific roles and responsibilities. For instance, the family socializes children, while education imparts necessary skills and knowledge, reinforcing societal norms. Functionalists argue that when all parts of society function harmoniously, social order is maintained; disruptions, such as crime or inequality, signal that something is amiss within the system. Thus, the theory emphasizes the importance of collective consensus and shared values in preserving social stability.
\\ \textit{Note:} The answer correctly identifies functionalist theory as a perspective that views society as an interconnected system where various social institutions work collaboratively to maintain order and stability, highlighting the roles of family, education, religion, and government in this process.
\end{enumerate}

\vfill
\noindent\textit{End of Answer Key}
\end{document}